\appendix
\section{Checklist de reproducibilidad (plantilla)}
\begin{itemize}[leftmargin=*]
  \item \textbf{Datos}: fuente, versión, licencias, anonimización.
  \item \textbf{Código}: repositorio, commit hash, instrucciones de ejecución.
  \item \textbf{Entorno}: SO, versión de compiladores, dependencias, semillas.
  \item \textbf{Procedimiento}: pasos exactos para replicar resultados.
  \item \textbf{Resultados}: tablas/figuras generadas automáticamente en \texttt{build/}.
\end{itemize}

\section{Diagramas de Arquitectura}

\subsection{Diagrama de Capas de Clean Architecture}
\begin{figure}[htbp]
    \centering
    \includegraphics[width=0.8\textwidth]{graphics/clean_architecture_diagram.pdf}
    \caption{Diagrama de capas de Clean Architecture implementado en Codexy}
    \label{fig:clean_architecture}
\end{figure}

\subsection{Diagrama de Flujo de Datos}
\begin{figure}[htbp]
    \centering
    \includegraphics[width=0.8\textwidth]{graphics/data_flow_diagram.pdf}
    \caption{Flujo de datos desde el escaneo QR hasta la actualización en tiempo real}
    \label{fig:data_flow}
\end{figure}

\section{Código Fuente Adicional}

\subsection{Configuración de SignalR en .NET 9}
\ifminted
\begin{minted}[linenos]{csharp}
// Startup.cs o Program.cs
builder.Services.AddSignalR();

app.MapHub<InventoryHub>("/inventoryHub");

// InventoryHub.cs
public class InventoryHub : Hub
{
    public async Task SendInventoryUpdate(string zoneId, InventoryUpdate update)
    {
        await Clients.Group(zoneId).SendAsync("ReceiveInventoryUpdate", update);
    }

    public async Task JoinZoneGroup(string zoneId)
    {
        await Groups.AddToGroupAsync(Context.ConnectionId, zoneId);
    }
}
\end{minted}
\else
\begin{lstlisting}[language=C#]
// Startup.cs o Program.cs
builder.Services.AddSignalR();

app.MapHub<InventoryHub>("/inventoryHub");

// InventoryHub.cs
public class InventoryHub : Hub
{
    public async Task SendInventoryUpdate(string zoneId, InventoryUpdate update)
    {
        await Clients.Group(zoneId).SendAsync("ReceiveInventoryUpdate", update);
    }

    public async Task JoinZoneGroup(string zoneId)
    {
        await Groups.AddToGroupAsync(Context.ConnectionId, zoneId);
    }
}
\end{lstlisting}
\fi

\subsection{Implementación de JWT en el Backend}
\ifminted
\begin{minted}[linenos]{csharp}
// AuthService.cs
public class AuthService : IAuthService
{
    private readonly IConfiguration _configuration;

    public AuthService(IConfiguration configuration)
    {
        _configuration = configuration;
    }

    public string GenerateJwtToken(User user)
    {
        var claims = new[]
        {
            new Claim(ClaimTypes.NameIdentifier, user.Id.ToString()),
            new Claim(ClaimTypes.Name, user.Username),
            new Claim(ClaimTypes.Role, user.Role)
        };

        var key = new SymmetricSecurityKey(Encoding.UTF8.GetBytes(_configuration["Jwt:Key"]));
        var creds = new SigningCredentials(key, SecurityAlgorithms.HmacSha256);

        var token = new JwtSecurityToken(
            issuer: _configuration["Jwt:Issuer"],
            audience: _configuration["Jwt:Audience"],
            claims: claims,
            expires: DateTime.Now.AddHours(8),
            signingCredentials: creds);

        return new JwtSecurityTokenHandler().WriteToken(token);
    }
}
\end{minted}
\else
\begin{lstlisting}[language=C#]
// AuthService.cs
public class AuthService : IAuthService
{
    private readonly IConfiguration _configuration;

    public AuthService(IConfiguration configuration)
    {
        _configuration = configuration;
    }

    public string GenerateJwtToken(User user)
    {
        var claims = new[]
        {
            new Claim(ClaimTypes.NameIdentifier, user.Id.ToString()),
            new Claim(ClaimTypes.Name, user.Username),
            new Claim(ClaimTypes.Role, user.Role)
        };

        var key = new SymmetricSecurityKey(Encoding.UTF8.GetBytes(_configuration["Jwt:Key"]));
        var creds = new SigningCredentials(key, SecurityAlgorithms.HmacSha256);

        var token = new JwtSecurityToken(
            issuer: _configuration["Jwt:Issuer"],
            audience: _configuration["Jwt:Audience"],
            claims: claims,
            expires: DateTime.Now.AddHours(8),
            signingCredentials: creds);

        return new JwtSecurityTokenHandler().WriteToken(token);
    }
}
\end{lstlisting}
\fi

\section{Métricas de Rendimiento Detalladas}

\begin{table}[htbp]
\centering
\caption{Métricas de rendimiento del sistema}
\label{tab:performance_metrics}
\begin{tabular}{|l|c|c|c|}
\hline
\textbf{Métrica} & \textbf{Valor Actual} & \textbf{Valor Objetivo} & \textbf{Status} \\
\hline
Tiempo de respuesta API & 150ms & <200ms & \checkmark \\
\hline
Tiempo de carga app móvil & 2.5s & <3s & \checkmark \\
\hline
Disponibilidad del sistema & 99.9\% & >99.5\% & \checkmark \\
\hline
Cobertura de pruebas & 85\% & >80\% & \checkmark \\
\hline
\end{tabular}
\end{table}

\section{Guía de Despliegue}

\subsection{Prerrequisitos}
\begin{itemize}
  \item .NET 9 SDK instalado
  \item Node.js 18+ para el frontend
  \item SQL Server 2019+
  \item Docker (opcional para contenedorización)
\end{itemize}

\subsection{Pasos de Despliegue}
\begin{enumerate}
  \item Clonar el repositorio: \texttt{git clone https://github.com/example/codexy.git}
  \item Configurar la base de datos: Ejecutar scripts en \texttt{database/setup.sql}
  \item Configurar variables de entorno en \texttt{appsettings.json}
  \item Construir el backend: \texttt{dotnet build}
  \item Ejecutar migraciones: \texttt{dotnet ef database update}
  \item Construir el frontend: \texttt{cd client && npm install && npm run build}
  \item Ejecutar la aplicación: \texttt{dotnet run}
\end{enumerate}

\section{Plan de Mantenimiento y Evolución}

\subsection{Mantenimiento Preventivo}
- Revisiones semanales de logs
- Actualizaciones mensuales de dependencias
- Auditorías trimestrales de seguridad

\subsection{Roadmap de Evolución}
- Q1 2025: Integración con IA para predicción de inventarios
- Q2 2025: Soporte para dispositivos IoT
- Q3 2025: Migración a microservicios si escala requiere
- Q4 2025: Internacionalización y soporte multiidioma

\section{Recursos Adicionales y Documentación}

\subsection{Repositorio y Documentación Técnica}
El código fuente completo de Codexy está disponible en GitHub bajo licencia MIT. La documentación técnica incluye:
- Guías de arquitectura detalladas
- Manuales de API con ejemplos de uso
- Tutoriales de despliegue para diferentes entornos
- Casos de uso y escenarios de prueba

\subsection{Referencias Bibliográficas Adicionales}
Además de las citadas en el texto principal:
- Fowler, M. (2003). Patterns of Enterprise Application Architecture. Addison-Wesley.
- Evans, E. (2003). Domain-Driven Design. Addison-Wesley.
- Martin, R. C. (2017). Clean Architecture. Prentice Hall.
- Beck, K. (2000). Extreme Programming Explained. Addison-Wesley.

\subsection{Herramientas y Tecnologías Utilizadas}
- Backend: .NET 9, Entity Framework Core, SignalR
- Frontend: Ionic 7, Angular 17, Capacitor
- Base de datos: SQL Server 2022
- Infraestructura: Azure DevOps, Docker, Kubernetes
- Testing: xUnit, Selenium, Cypress
- Monitoreo: Application Insights, Prometheus

\subsection{Equipo de Desarrollo}
- Arquitecto de Software: [Nombre]
- Desarrolladores Backend: [Nombres]
- Desarrolladores Frontend: [Nombres]
- QA Engineers: [Nombres]
- Product Owner: [Nombre]
- Scrum Master: [Nombre]

Este apéndice proporciona recursos completos para replicar, extender o mantener el sistema Codexy.

\section{Glosario de Términos}

- **Clean Architecture**: Enfoque de diseño que separa el software en capas concéntricas, promoviendo independencia de frameworks y facilidad de testing.

- **SignalR**: Biblioteca para comunicación bidireccional en tiempo real entre cliente y servidor.

- **JWT**: JSON Web Token, estándar para autenticación segura y stateless.

- **Entity Framework Core**: ORM para .NET que facilita el acceso a bases de datos.

- **Ionic**: Framework para desarrollo de aplicaciones móviles híbridas usando web technologies.

- **DevOps**: Cultura y prácticas que integran desarrollo de software con operaciones IT.

\section{Cronograma del Proyecto}

\begin{table}[htbp]
\centering
\caption{Cronograma de desarrollo de Codexy}
\label{tab:timeline}
\begin{tabular}{|l|c|c|}
\hline
\textbf{Fase} & \textbf{Duración} & \textbf{Hitos Principales} \\
\hline
Planificación & 2 semanas & Requerimientos, arquitectura inicial \\
\hline
Desarrollo Backend & 8 semanas & API, base de datos, lógica de negocio \\
\hline
Desarrollo Frontend & 6 semanas & App móvil, interfaz web \\
\hline
Integración y Testing & 4 semanas & Pruebas end-to-end, optimizaciones \\
\hline
Despliegue y Capacitación & 2 semanas & Producción, entrenamiento usuarios \\
\hline
\end{tabular}
\end{table}

\section{Lecciones Aprendidas Detalladas}

Durante el desarrollo, aprendimos valiosas lecciones que pueden beneficiar futuros proyectos:

1. **Importancia de la Comunicación**: Reuniones diarias y documentación clara previnieron malentendidos.

2. **Valor de la Automatización**: Herramientas CI/CD redujeron errores humanos y aceleraron entregas.

3. **Necesidad de Flexibilidad**: Cambios en requerimientos son inevitables; la agilidad es clave.

4. **Enfoque en Calidad**: Testing temprano previene problemas costosos en producción.

5. **Colaboración Interdisciplinaria**: Equipos diversos producen soluciones más robustas.